\subsection{共轭扩张与共轭元素}

定义 1.——设 E 与 F 是包含在 Ω 中的 K 的两个扩张。若存在 Ω 的一个 K-自同构 u 使得 \(u(E) = F\) ，则称 E 与 F （在 Ω 中）共轭。Ω 的两个元素 x 与 y 称为在 K 上共轭，如果存在 Ω 的一个 K-自同构 u 使得 \(u(x) = y\) .

设 E 与 F 是包含在 \(\Omega\) 中的 K 的两个扩张。由命题 1，E 与 F 在 K 上共轭当且仅当它们是 K 的同构扩张。特别地，当存在 \(\Omega\) 的两个子集 A 与 B 使得 \(\mathrm{E} = \mathrm{K}(\mathrm{A})\) 且 \(\mathrm{F} = \mathrm{K}(\mathrm{B})\) ，并且存在 \(\Omega\) 的一个 K-自同构 u 使得 \(u(\mathrm{A}) = \mathrm{B}\) 时，即是如此。

关系 « x 与 y 在 K 上共轭 » 是 \(\Omega\) 中的一个等价关系；其等价类称为 \(\Omega\) 中的共轭类；它们是 \(\Omega\) 的全体 K-自同构构成的群在 \(\Omega\) 中的轨道。

命题 2.——设 \(x\) 与 \(y\) 是 \(\Omega\) 的两个元素。则下列条件等价：

\begin{enumerate}
\def\labelenumi{\alph{enumi})}
\item
  \(x\) 与 \(y\) 在 \(\mathbf{K}\) 上共轭。
\item
  存在一个 K-同构 \(v\) （\(\mathbf{K}(x)\) 到 \(\mathbf{K}(y)\) 上的），使得 \(v(x) = y\) 。
\item
  \(x\) 与 \(y\) 在 \(\mathbf{K}\) 上有相同的极小多项式。
\end{enumerate}

首先设 x 与 y 在 K 上共轭；设 u 是 \(\Omega\) 的一个 K-自同构，使得 \(u(x) = y\) ，并设 f 是 x 在 K 上的极小多项式。我们有

\[
f (y) = f (u (x)) = u (f (x)) = 0,
\]

且 \(f\) 是 \(\mathbf{K}[\mathbf{X}]\) 中的首一不可约多项式；因此（V，第 16 页，定理 1，c)），\(f\) 是 \(y\) 在 \(\mathbf{K}\) 上的极小多项式。于是 \(a)\) 蕴含 \(c)\) 。

现在设 x 与 y 在 K 上有相同的极小多项式 f。存在域 \(\mathrm{K}[\mathrm{X}]/(f)\) 到 \(\mathrm{K}(x)\) （相应地，到 \(\mathrm{K}(y)\) ）上的一个 K-同构，它把 X 模 (f) 的剩余类映到 x（相应地，y）（V，第 16 页，定理 1，b)）；因此存在 \(\mathrm{K}(x)\) 到 \(\mathrm{K}(y)\) 上的一个 K-同构 v，使得 \(v(x)=y\) 。故 c) 蕴含 b)。

最后，在假设 \(b\) ) 下，命题 1 蕴含存在一个 K-自同构 \(u\) （\(\Omega\) 的），它延拓 \(v\) ；于是我们有 \(u(x) = y\) ，从而 \(x\) 与 \(y\) 在 \(K\) 上共轭，故 \(b\) ) 蕴含 \(a\) )。

推论 1.——设 x 是 Ω 中的元素，在 K 上次数为 n，并设 f 是 x 在 K 上的极小多项式。x 在 K 上的共轭元是 f 在 Ω 中的根，其个数至多等于 n。

推论 2.——设 x 是 Ω 中的元素，在 K 上次数为 n。x 在 K 上可分当且仅当 x 在 Ω 中有 n 个共轭元；当此情形成立时，x 在 K 上的所有共轭元都在 K 上可分。

设 \(f\) 是 \(x\) 在 \(K\) 上的极小多项式；它的根是 \(x\) 在 \(K\) 上的共轭元，且这些根中的每一个都以 \(f\) 为其在 \(K\) 上的极小多项式。而 \(x\) 在 \(K\) 上可分当且仅当多项式 \(f\) 是可分的（V，第 39 页，命题 5），即当且仅当 \(f\) 有 \(n\) 个互不相同的根（在 \(\Omega\) 中）。推论 2 由此得证。

推论 3.——设 G 是 \(\Omega\) 的 K-自同构的群。G 在 \(\Omega\) 中的不变量集是相对 \(p\) -根式闭包（K 在 \(\Omega\) 中的）（V，第 25 页）。

换言之，一个元素 \(x\) （\(\Omega\) 的）是 \(p\) -根式的（在 \(\mathbf{K}\) 上）当且仅当它除自身外没有别的共轭元。

设 \(p\) 为特征指数，并设 \(x \in \Omega\) 。由 V 第 44 页命题 13，存在整数 \(m \geqslant 0\) 使得 \(y = x^{p^m}\) 在 K 上是代数且可分的。而 \(x\) 在 G 下不变当且仅当 \(y\) 在 G 下不变；由推论 2，这当且仅当 \(y\) 在 K 上次数为 1，即 \(y \in K\) 。由此推论立即可得。

